% !TeX program = pdflatex
\documentclass[11pt,a4paper]{amsart}

\newcommand{\notelanguage}{finnish}
% Shared preamble for course notes and exercise sets.

\usepackage[T1]{fontenc}
\usepackage{lmodern}
\providecommand{\notelanguage}{english}
\usepackage[finnish,english]{babel}
\AtBeginDocument{\expandafter\selectlanguage\expandafter{\notelanguage}}

\newcommand{\notesfinnishlanguage}{finnish}
\ifx\notelanguage\notesfinnishlanguage
\newcommand{\notestheoremname}{Lause}
\newcommand{\noteslemmaname}{Lemma}
\newcommand{\notespropositionname}{Propositio}
\newcommand{\notescorollaryname}{Seurauslause}
\newcommand{\notesdefinitionname}{Määritelmä}
\newcommand{\notesexamplename}{Esimerkki}
\newcommand{\notesproblemname}{Tehtävä}
\newcommand{\notesexercisename}{Tehtävä}
\newcommand{\notesremarkname}{Huomautus}
\newcommand{\notessolutionname}{Ratkaisu}
\else
\newcommand{\notestheoremname}{Theorem}
\newcommand{\noteslemmaname}{Lemma}
\newcommand{\notespropositionname}{Proposition}
\newcommand{\notescorollaryname}{Corollary}
\newcommand{\notesdefinitionname}{Definition}
\newcommand{\notesexamplename}{Example}
\newcommand{\notesproblemname}{Problem}
\newcommand{\notesexercisename}{Exercise}
\newcommand{\notesremarkname}{Remark}
\newcommand{\notessolutionname}{Solution}
\fi

\usepackage[a4paper,margin=30mm]{geometry}
\usepackage{microtype}
\usepackage{graphicx}
\usepackage{enumitem}

\newcommand{\exerciseheading}[2]{%
  \par\noindent
  {\large\bfseries #1\par}
  \vspace{0.1em}
  \noindent{\large #2\par}
  \vspace{1.5em}
}

\usepackage{amsmath}
\usepackage{amssymb}
\usepackage{amsthm}
\usepackage{mathtools}
\usepackage{aliascnt}

\usepackage[hidelinks,pdfusetitle]{hyperref}

\numberwithin{equation}{section}
\setlist{itemsep=0.25em,topsep=0.5em}

\theoremstyle{plain}
\newtheorem{theorem}{\notestheoremname}[section]
\newaliascnt{lemma}{theorem}
\newtheorem{lemma}[lemma]{\noteslemmaname}
\aliascntresetthe{lemma}
\newaliascnt{proposition}{theorem}
\newtheorem{proposition}[proposition]{\notespropositionname}
\aliascntresetthe{proposition}
\newaliascnt{corollary}{theorem}
\newtheorem{corollary}[corollary]{\notescorollaryname}
\aliascntresetthe{corollary}

\theoremstyle{definition}
\newaliascnt{definition}{theorem}
\newtheorem{definition}[definition]{\notesdefinitionname}
\aliascntresetthe{definition}
\newaliascnt{example}{theorem}
\newtheorem{example}[example]{\notesexamplename}
\aliascntresetthe{example}
\newaliascnt{problem}{theorem}
\newtheorem{problem}[problem]{\notesproblemname}
\aliascntresetthe{problem}
\newtheorem{exercise}{\notesexercisename}

\theoremstyle{remark}
\newaliascnt{remark}{theorem}
\newtheorem{remark}[remark]{\notesremarkname}
\aliascntresetthe{remark}

\newenvironment{solution}
{
\begin{proof}[\notessolutionname]}
  {
\end{proof}}

\usepackage[nameinlink,noabbrev]{cleveref}

\crefname{theorem}{theorem}{theorems}
\crefname{lemma}{lemma}{lemmas}
\crefname{proposition}{proposition}{propositions}
\crefname{corollary}{corollary}{corollaries}
\crefname{definition}{definition}{definitions}
\crefname{example}{example}{examples}
\crefname{problem}{problem}{problems}
\crefname{exercise}{exercise}{exercises}
\crefname{remark}{remark}{remarks}

\ifx\notelanguage\notesfinnishlanguage
\crefname{theorem}{lause}{lauseet}
\crefname{lemma}{lemma}{lemmat}
\crefname{proposition}{propositio}{propositiot}
\crefname{corollary}{seurauslause}{seurauslauseet}
\crefname{definition}{määritelmä}{määritelmät}
\crefname{example}{esimerkki}{esimerkit}
\crefname{problem}{tehtävä}{tehtävät}
\crefname{exercise}{tehtävä}{tehtävät}
\crefname{remark}{huomautus}{huomautukset}
\fi

\newcommand{\NN}{\mathbb{N}}
\newcommand{\NNone}{\mathbb{N}_1}
\newcommand{\ZZ}{\mathbb{Z}}
\newcommand{\QQ}{\mathbb{Q}}
\newcommand{\RR}{\mathbb{R}}
\newcommand{\CC}{\mathbb{C}}
\newcommand{\PP}{\mathbb{P}}

\DeclareMathOperator{\im}{im}
\DeclareMathOperator{\rank}{rank}
\DeclarePairedDelimiter{\abs}{\lvert}{\rvert}
\DeclarePairedDelimiter{\norm}{\lVert}{\rVert}
\newcommand{\set}[1]{\left\lbrace#1\right\rbrace}
\newcommand{\maxset}[1]{\max\set{#1}}
\newcommand{\minset}[1]{\min\set{#1}}
\newcommand{\powerset}[1]{\mathcal{P}\!\left(#1\right)}


\title{Analyysi I\\Ryhmätyö 6}
\author{}
\date{}

\begin{document}

\exerciseheading{Analyysi I}{Ryhmätyö 6}

\begin{exercise}
  \leavevmode\par
  \begin{enumerate}[label=(\alph*), nosep]
    \item Anna esimerkki funktiosta \(f\colon \RR \to \RR\), jolla ei ole
      oikeanpuoleista raja-arvoa origossa.
    \item Voiko jonoille määritellä oikeanpuoleinen ja vasemmanpuoleinen
      raja-arvo?
    \item Määrittele analogian avulla raja-arvot
      \(\lim_{x \to x_0^{\pm}} f(x) = \pm\infty\).
    \item Anna esimerkki funktiosta \(f\), jolle
      \(\lim_{x \to 1^-} f(x) = \infty\).
  \end{enumerate}
\end{exercise}

\begin{solution}
  \leavevmode\par
  \begin{enumerate}[label=(\alph*), nosep]
    \item Olkoon \(f \colon \RR \to \RR\),
      \[
        f(x) =
        \begin{cases}
          \sin(1/x), & x > 0, \\
          0, & x \leq 0.
        \end{cases}
      \]
      Jonot \(x_n = 1/(\pi/2 + 2\pi n)\) ja
      \(y_n = 1/(3\pi/2 + 2\pi n)\), \(n \geq 1\), lähestyvät nollaa
      oikealta, mutta \(f(x_n) = 1\) ja \(f(y_n) = -1\) kaikilla \(n\).
      Täten funktiolla \(f\) ei ole oikeanpuoleista raja-arvoa origossa.
    \item Tavalliselle jonon raja-arvolle ei määritellä erikseen oikean- ja
      vasemmanpuoleista raja-arvoa kuten funktion raja-arvolle äärellisessä
      pisteessä. Jonon raja-arvossa indeksi \(n\) lähestyy ääretöntä, eikä
      mitään äärellistä pistettä kahdelta eri puolelta. Funktion
      toispuoleisissa raja-arvoissa puoli tarkoittaa argumentin
      lähestymissuuntaa, ei funktion arvojen sijaintia raja-arvoon nähden.
      Jonon termit voivat lähestyä raja-arvoa yläpuolelta, alapuolelta tai
      vuorotellen molemmilta puolilta, kuten jonot \(1/n\), \(-1/n\) ja
      \((-1)^n/n\) lähestyvät nollaa.
    \item Olkoon \(f \colon A \to \RR\).

      Oletetaan, että on olemassa \(r > 0\) siten, että
      \((x_0 - r, x_0) \subset A\). Tällöin määritellään
      \(\lim_{x \to x_0^-} f(x) = +\infty\) jos ja vain jos jokaista
      \(M \in \RR\) kohti on olemassa sellainen \(0 < \delta \leq r\), että
      \(f(x) > M\) kaikilla \(x \in (x_0 - \delta, x_0)\).

      Oletetaan, että on olemassa \(r > 0\) siten, että
      \((x_0 - r, x_0) \subset A\). Tällöin määritellään
      \(\lim_{x \to x_0^-} f(x) = -\infty\) jos ja vain jos jokaista
      \(m \in \RR\) kohti on olemassa sellainen \(0 < \delta \leq r\), että
      \(f(x) < m\) kaikilla \(x \in (x_0 - \delta, x_0)\).

      Oletetaan, että on olemassa \(r > 0\) siten, että
      \((x_0, x_0 + r) \subset A\). Tällöin määritellään
      \(\lim_{x \to x_0^+} f(x) = +\infty\) jos ja vain jos jokaista
      \(M \in \RR\) kohti on olemassa sellainen \(0 < \delta \leq r\), että
      \(f(x) > M\) kaikilla \(x \in (x_0, x_0 + \delta)\).

      Oletetaan, että on olemassa \(r > 0\) siten, että
      \((x_0, x_0 + r) \subset A\). Tällöin määritellään
      \(\lim_{x \to x_0^+} f(x) = -\infty\) jos ja vain jos jokaista
      \(m \in \RR\) kohti on olemassa sellainen \(0 < \delta \leq r\), että
      \(f(x) < m\) kaikilla \(x \in (x_0, x_0 + \delta)\).

    \item Olkoon \(f \colon \RR \setminus \set{1} \to \RR\), \(f(x) =
      1/\abs{x - 1}\). Olkoon \(M \in \RR\). Jos \(M \leq 0\), valitaan
      \(\delta = 1\). Tällöin \(f(x) > 0 \geq M\), kun \(x \in (1 -
      \delta, 1)\).
      Oletetaan sitten, että \(M > 0\). Valitaan \(\delta = \minset{1, 1/M}\).
      Tällöin
      \[
        f(x) = \frac{1}{\abs{x - 1}} = \frac{1}{1 - x} >
        \frac{1}{\delta} \geq M,
      \]
      kun \(1 - \delta < x < 1\).
  \end{enumerate}
\end{solution}

\begin{exercise}
  Tutustu kirjan Lauseeseen 3.2.3. Voiko lauseen todistusta muokata niin,
  että se toimii tapauksessa \(m = \infty\)?
\end{exercise}

\begin{solution}
  Kyllä. Korvataan todistuksessa ehto \(\abs{f(x) - m} < \varepsilon\)
  ehdolla \(f(x) > M\), missä \(M \in \RR\) on mielivaltainen.
  Lauseen oletuksilla pätee siis
  \[
    \lim_{x \to x_0} f(x) = +\infty
    \quad\Longleftrightarrow\quad
    \lim_{x \to x_0^-} f(x) = \lim_{x \to x_0^+} f(x) = +\infty.
  \]

  Oletetaan ensin, että \(\lim_{x \to x_0} f(x) = +\infty\), ja olkoon
  \(M \in \RR\). Tällöin on olemassa \(0 < \delta \leq r\) siten, että
  \(f(x) > M\), kun \(0 < \abs{x - x_0} < \delta\). Ehto pätee siis
  sekä silloin, kun \(x_0 - \delta < x < x_0\), että silloin, kun
  \(x_0 < x < x_0 + \delta\). Näin ollen molemmat toispuoleiset
  raja-arvot ovat \(+\infty\).

  Oletetaan sitten, että molemmat toispuoleiset raja-arvot ovat
  \(+\infty\), ja olkoon \(M \in \RR\). Tällöin on olemassa
  \(\delta_-, \delta_+ > 0\) siten, että
  \[
    f(x) > M \quad\text{kun } x_0 - \delta_- < x < x_0,
  \]
  ja
  \[
    f(x) > M \quad\text{kun } x_0 < x < x_0 + \delta_+.
  \]
  Valitaan \(\delta = \min\{r, \delta_-, \delta_+\} > 0\).
  Jos \(0 < \abs{x - x_0} < \delta\), niin \(x\) kuuluu jompaankumpaan
  yllä olevista väleistä, joten \(f(x) > M\). Koska \(M\) oli
  mielivaltainen, saadaan \(\lim_{x \to x_0} f(x) = +\infty\).
\end{solution}

\end{document}
